\begin{problem}[33届物理竞赛决赛]{64}
    一、山西大同某煤矿相对于秦皇岛的高度为 $h_{c}$ 。质量为 $m_{t}$ 的火车载有质量为 $m_{c}$ 的煤，从大同沿大秦线铁路行驶路程 $l$ 后到达秦皇岛，卸载后空车返回。从大同到秦皇岛的过程中，火车和煤总势能的一部分克服铁轨和空气阻力做功，其余部分由发电机转换成电能，平均转换效率为 $\eta_{1}$ ，电能被全部储存于蓄电池中以用于返程。空车在返程中由储存的电能驱动电动机克服重力和阻力做功，存储电能转化为对外做功的平均转换效率为 $\eta_{2}$ 。假设大秦线轨道上火车平均每运行单位距离克服阻力需要做的功与运行时（火车或火车和煤）总重量成正比，比例系数为常数 $\mu$ ，火车由大同出发时携带的电能为零。
    \begin{subproblem}
        若空车返回大同时还有剩余的电能，求该电能 $E$。
    \end{subproblem}
    \begin{subproblem}
        问火车至少装载质量为多少的煤，才能在不另外提供能量的条件下刚好返回大同？
    \end{subproblem}
    \begin{subproblem}
        已知火车在从大同到秦皇岛的铁轨上运行的平均速率为 $v$ ，请给出发电机的平均输出功率 $P$ 与题给的其它物理量的关系。
    \end{subproblem}
\end{problem}